\subsection{局部池化与非线性变换}
\label{l11:subsec:pooling}

池化把一个局部窗口内的多个数值汇总为一个输出。窗口沿空间轴移动，输入通道之间保持分离，因此常规空间池化保留批次大小和通道数量，同时改变每个通道的空间尺寸。常见汇总函数包括平均值、最大值、$L^2$ 范数和加权平均。

设一个窗口包含 $N^2$ 个值，按任意固定顺序记为 $z_1,\ldots,z_{N^2}$。平均池化输出 $N^{-2}\sum_i z_i$，最大池化输出 $\max_i z_i$，$L^2$ 池化输出 $(\sum_i z_i^2)^{1/2}$；带非负权重 $a_i$ 且 $\sum_i a_i=1$ 的加权平均输出 $\sum_i a_i z_i$。这些选择保留的局部信息不同：平均值反映整体响应，最大值保留最强响应，$L^2$ 范数反映窗口中响应的合成幅度。若将平方和先除以 $N^2$ 再开平方，得到的是均方根形式，需要与未归一化的 $L^2$ 范数区分。

对卷积输出 $Y$ 使用 $N\times N$ 窗口、步幅为 $N$、无填充的不重叠平均池化，空间输出尺寸为 $H_{\mathrm{s}}=\left\lfloor H_{\mathrm{o}}/N\right\rfloor$、$W_{\mathrm{s}}=\left\lfloor W_{\mathrm{o}}/N\right\rfloor$。池化输出坐标 $(x,y)$ 所需窗口的左上角是 $(Nx,Ny)$，因此窗口内元素为 $Y_{b,m,Nx+p,Ny+q}$，其中 $0\le p,q<N$。若尺寸无法整除 $N$，最右侧或最下侧不足一个完整窗口的元素不参与这一约定下的池化。

进一步对平均值加上输出通道偏置 $\beta_m$，再施加 S 形激活函数（sigmoid），得到
\begin{equation}
\begin{aligned}
  S_{b,m,x,y}
  &=\sigma\!\left(
     \frac{1}{N^2}\sum_{p=0}^{N-1}\sum_{q=0}^{N-1}
     Y_{b,m,Nx+p,Ny+q}+\beta_m
    \right),\\
  \sigma(z)&=\frac{1}{1+\exp(-z)},\qquad
  S\in\mathbb{R}^{B\times M\times H_{\mathrm{s}}\times W_{\mathrm{s}}}.
\end{aligned}
\label{l11:eq:pool_sigmoid}
\end{equation}
偏置 $\beta_m$ 在所有图像和同一输出通道的全部空间位置上共享，所以这个偏置向量只包含 $M$ 个数。平均池化本身没有需要学习的权重；附加偏置以后，这个组合模块具有 $M$ 个可学习参数。

式\eqref{l11:eq:pool_sigmoid}的运算顺序是先求平均、再加偏置、最后激活。非线性函数通常不能与平均运算交换。例如，窗口值为 $\{0,0,2,2\}$ 且偏置为 $0$ 时，先平均再激活得到 $\sigma(1)\approx0.7311$；先激活再平均则得到 $[\sigma(0)+\sigma(2)]/2\approx0.6904$。两者表示不同的层结构。固定通道偏置与平均可以交换，是因为 $\operatorname{avg}(Y+\beta_m)=\operatorname{avg}(Y)+\beta_m$；这项线性关系不能用来交换后面的非线性激活。

激活函数还可以使用整流线性单元（Rectified Linear Unit，ReLU），定义为 $\phi(z)=\max(0,z)$。它与 S 形函数都逐元素作用，不改变张量形状。具体实现应按照网络给定的函数和顺序执行，不能因为两个模块都称为激活层就互相替换。

\subsection{池化的数值解释与独立性}

对式\eqref{l11:eq:channel_example_output}中的整个 $2\times2$ 输出做一次平均池化，得到 $(51+50+60+48)/4=52.25$。若偏置取 $-52$，则激活输入为 $0.25$，最终输出为 $\sigma(0.25)\approx0.5622$。若只做最大池化，输出为 $60$；若做未归一化的 $L^2$ 池化，输出为 $\sqrt{51^2+50^2+60^2+48^2}=\sqrt{11\,005}$。同一个窗口的不同汇总规则会产生不同结果，不能只根据输出尺寸识别池化操作。

池化可以降低对局部位置变化的敏感性，但这种效果取决于窗口边界与汇总方式。以一维最大池化作为简化例子，窗口宽度和步幅均为 $2$：输入 $[9,0,0,0]$ 和 $[0,9,0,0]$ 都得到 $[9,0]$；若响应继续右移一个位置，输入变为 $[0,0,9,0]$，输出就变为 $[0,9]$。因此，在一个窗口内部移动的特征可能给出相同结果，跨过窗口边界后则可能改变输出。空间下采样本身也不保证尺度不变性。

从并行计算角度看，$(b,m,x,y)$ 不同的池化输出可以独立计算。一个线程可以读取所需的 $N^2$ 个值，在寄存器中累加，除以 $N^2$，再执行偏置与激活，最后写回一个输出。这个方案与卷积中的单输出线程具有相同的组织原则：输出由线程唯一负责，局部窗口内的求和由该线程完成。

卷积与池化之间还有一项可利用的数据复用。如果池化所需的卷积结果已由同一线程块产生，并且在写出前仍位于寄存器或共享内存中，就可以直接完成局部汇总，减少中间张量的全局写回和随后重读。合并实现时，必须保留式\eqref{l11:eq:pool_sigmoid}的准确顺序，并确保每个池化窗口所需的全部卷积结果已经可用。窗口跨越线程块边界时，块内同步不能使另一个块的中间结果自动变得可用，需要重新安排分工或保留独立执行阶段。
